\documentclass[10pt,a4paper]{article}

% ------------------------------------------------------------
% PACKAGES
% ------------------------------------------------------------
\usepackage[a4paper,margin=1.25cm]{geometry}
\usepackage{amsmath,amssymb}
\usepackage{xcolor}
\usepackage{listings}
\usepackage{tcolorbox}
\usepackage{multicol}
\usepackage{enumitem}
\usepackage{titlesec}
\usepackage{hyperref}
\usepackage{microtype}

% ------------------------------------------------------------
% COLORS
% ------------------------------------------------------------
\definecolor{navy}{RGB}{25,55,95}
\definecolor{lightblue}{RGB}{235,242,250}
\definecolor{lightgray}{RGB}{245,245,245}
\definecolor{darkgray}{RGB}{70,70,70}
\definecolor{codegray}{RGB}{245,245,245}

% ------------------------------------------------------------
% PAGE STYLE
% ------------------------------------------------------------
\pagestyle{empty}

\setlength{\parindent}{0pt}
\setlength{\parskip}{2pt}

% ------------------------------------------------------------
% SECTION STYLE
% ------------------------------------------------------------
\titleformat{\section}
  {\large\bfseries\color{navy}}
  {}
  {0pt}
  {}

\titlespacing*{\section}{0pt}{5pt}{3pt}

% ------------------------------------------------------------
% PYTHON CODE STYLE
% ------------------------------------------------------------
\lstdefinestyle{pythonstyle}{
    language=Python,
    basicstyle=\ttfamily\scriptsize,
    keywordstyle=\color{navy}\bfseries,
    commentstyle=\color{darkgray}\itshape,
    stringstyle=\color{red!60!black},
    showstringspaces=false,
    keepspaces=true,
    columns=fullflexible,
    frame=single,
    rulecolor=\color{black!15},
    backgroundcolor=\color{codegray},
    breaklines=true,
    tabsize=4,
    xleftmargin=3pt,
    xrightmargin=3pt,
    aboveskip=2pt,
    belowskip=2pt
}

% ------------------------------------------------------------
% CONCEPT BOX
% ------------------------------------------------------------
\newtcolorbox{card}{
    colback=white,
    colframe=black!12,
    boxrule=0.4pt,
    arc=2pt,
    left=4pt,
    right=4pt,
    top=3pt,
    bottom=3pt,
    before skip=2pt,
    after skip=2pt
}
\allowdisplaybreaks
% ------------------------------------------------------------
% DOCUMENT
% ------------------------------------------------------------
\begin{document}

% ============================================================
% TITLE
% ============================================================

\begin{center}
    {\LARGE\bfseries\color{navy}
    Scientific Python --- Quick Reference}\\[3pt]
    {\small
    Python, NumPy, Matplotlib, SciPy, and SymPy for physics and scientific computing --- with everything the Statistics book uses}
\end{center}

\vspace{3pt}

% ============================================================
% PAGE 1
% ============================================================

\begin{multicols}{2}
\raggedcolumns

% ------------------------------------------------------------
\section*{1. Python Basics}
% ------------------------------------------------------------

\begin{card}
\textbf{Variables} --- A name refers to an object.

\begin{lstlisting}[style=pythonstyle]
x = 3
y = 2.5
z = x + y

print(x)
print(z)
\end{lstlisting}

Output:
\texttt{3}\\
\texttt{5.5}
\end{card}

\begin{card}
\textbf{Basic types}

\begin{lstlisting}[style=pythonstyle]
a = 3
b = 3.14
c = 2 + 3j
d = "hello"
e = True
f = None

print(type(a))
print(type(c))
print(type(d))
print(type(e))
print(type(f))
\end{lstlisting}
\end{card}

\begin{card}
\textbf{\texttt{print()}} --- Display values.

\begin{lstlisting}[style=pythonstyle]
x = 3.14
print(x)
print("x =", x)
\end{lstlisting}

Output:
\texttt{3.14}\\
\texttt{x = 3.14}
\end{card}

\begin{card}
\textbf{\texttt{type()}} --- Find the type of an object.

\begin{lstlisting}[style=pythonstyle]
x = 3
print(type(x))
\end{lstlisting}

Output:
\texttt{<class 'int'>}
\end{card}

\begin{card}
\textbf{\texttt{isinstance()}} --- Test whether an object has a given type.

\begin{lstlisting}[style=pythonstyle]
x = 3.0

print(isinstance(x, float))
print(isinstance(x, int))
\end{lstlisting}

Output:
\texttt{True}\\
\texttt{False}
\end{card}

% ------------------------------------------------------------
\section*{2. Operators}
% ------------------------------------------------------------

\begin{card}
\textbf{Arithmetic}

\begin{lstlisting}[style=pythonstyle]
print(7 + 2)
print(7 - 2)
print(7 * 2)
print(7 / 2)
print(7 // 2)
print(7 % 2)
print(7 ** 2)
\end{lstlisting}

Output:
\texttt{9}\\
\texttt{5}\\
\texttt{14}\\
\texttt{3.5}\\
\texttt{3}\\
\texttt{1}\\
\texttt{49}
\end{card}

\begin{card}
\textbf{Comparison}

\begin{lstlisting}[style=pythonstyle]
x = 3

print(x == 3)
print(x != 3)
print(x > 2)
print(x >= 3)
print(x < 5)
print(x <= 2)
\end{lstlisting}

Output:
\texttt{True}\\
\texttt{False}\\
\texttt{True}\\
\texttt{True}\\
\texttt{True}\\
\texttt{False}
\end{card}

\begin{card}
\textbf{Logical operators}

\begin{lstlisting}[style=pythonstyle]
x = 3

print(x > 0 and x < 5)
print(x < 0 or x > 2)
print(not (x > 0))
\end{lstlisting}

Output:
\texttt{True}\\
\texttt{True}\\
\texttt{False}
\end{card}

\begin{card}
\textbf{\texttt{in}} --- Test membership.

\begin{lstlisting}[style=pythonstyle]
x = [1, 2, 3]

print(2 in x)
print(5 in x)
\end{lstlisting}

Output:
\texttt{True}\\
\texttt{False}
\end{card}

% ------------------------------------------------------------
\section*{3. Strings}
% ------------------------------------------------------------

\begin{card}
\textbf{Strings} --- Text enclosed by quotes.

\begin{lstlisting}[style=pythonstyle]
name = "Chandra"

print(name)
print(name[0])
print(name[-1])
\end{lstlisting}

Output:
\texttt{Chandra}\\
\texttt{C}\\
\texttt{a}
\end{card}

\begin{card}
\textbf{String slicing}

\begin{lstlisting}[style=pythonstyle]
s = "abcdef"

print(s[1:4])
print(s[:3])
print(s[3:])
print(s[::-1])
\end{lstlisting}

Output:
\texttt{bcd}\\
\texttt{abc}\\
\texttt{def}\\
\texttt{fedcba}
\end{card}

\begin{card}
\textbf{Useful string methods}

\begin{lstlisting}[style=pythonstyle]
s = "  Hello Python  "

print(s.strip())
print(s.lower())
print(s.upper())
print(s.replace("Python", "World"))
\end{lstlisting}
\end{card}

\begin{card}
\textbf{f-strings} --- Insert values directly into text.

\begin{lstlisting}[style=pythonstyle]
x = 3.14159

print(f"x = {x}")
print(f"x = {x:.2f}")
\end{lstlisting}

Output:
\texttt{x = 3.14159}\\
\texttt{x = 3.14}
\end{card}

% ------------------------------------------------------------
\section*{4. Conditions and Loops}
% ------------------------------------------------------------

\begin{card}
\textbf{\texttt{if / elif / else}} --- Conditional logic.

\begin{lstlisting}[style=pythonstyle]
x = -2

if x > 0:
    print("positive")
elif x < 0:
    print("negative")
else:
    print("zero")
\end{lstlisting}

Output:
\texttt{negative}
\end{card}

\begin{card}
\textbf{\texttt{for}} --- Iterate over elements.

\begin{lstlisting}[style=pythonstyle]
for i in range(5):
    print(i)
\end{lstlisting}

Output:
\texttt{0 1 2 3 4}
\end{card}

\begin{card}
\textbf{\texttt{range()}} --- Generate integer sequences.

\begin{lstlisting}[style=pythonstyle]
print(list(range(5)))
print(list(range(2, 7)))
print(list(range(0, 10, 2)))
\end{lstlisting}

Output:
\texttt{[0, 1, 2, 3, 4]}\\
\texttt{[2, 3, 4, 5, 6]}\\
\texttt{[0, 2, 4, 6, 8]}
\end{card}

\begin{card}
\textbf{\texttt{while}} --- Repeat while a condition is true.

\begin{lstlisting}[style=pythonstyle]
i = 0

while i < 3:
    print(i)
    i += 1
\end{lstlisting}

Output:
\texttt{0 1 2}
\end{card}

\begin{card}
\textbf{\texttt{break}} --- Exit a loop immediately.

\begin{lstlisting}[style=pythonstyle]
for i in range(10):
    if i == 3:
        break
    print(i)
\end{lstlisting}

Output:
\texttt{0 1 2}
\end{card}

\begin{card}
\textbf{\texttt{continue}} --- Skip the current iteration.

\begin{lstlisting}[style=pythonstyle]
for i in range(5):
    if i == 2:
        continue
    print(i)
\end{lstlisting}

Output:
\texttt{0 1 3 4}
\end{card}

\begin{card}
\textbf{\texttt{pass}} --- Do nothing; useful as a placeholder.

\begin{lstlisting}[style=pythonstyle]
for i in range(3):
    pass

print("done")
\end{lstlisting}

Output:
\texttt{done}
\end{card}

% ------------------------------------------------------------
\section*{5. Built-in Functions}
% ------------------------------------------------------------

\begin{card}
\textbf{\texttt{len()}} --- Number of elements.

\begin{lstlisting}[style=pythonstyle]
x = [2, 4, 6]

print(len(x))
\end{lstlisting}

Output: \texttt{3}
\end{card}

\begin{card}
\textbf{\texttt{sum()}} --- Sum elements.

\begin{lstlisting}[style=pythonstyle]
x = [1, 2, 3]

print(sum(x))
\end{lstlisting}

Output: \texttt{6}
\end{card}

\begin{card}
\textbf{\texttt{min()} / \texttt{max()}} --- Smallest/largest value.

\begin{lstlisting}[style=pythonstyle]
x = [4, 2, 9, 1]

print(min(x))
print(max(x))
\end{lstlisting}

Output:
\texttt{1}\\
\texttt{9}
\end{card}

\begin{card}
\textbf{\texttt{abs()}} --- Absolute value.

\begin{lstlisting}[style=pythonstyle]
print(abs(-3.5))
\end{lstlisting}

Output: \texttt{3.5}
\end{card}

\begin{card}
\textbf{\texttt{sorted()}} --- Return a sorted copy.

\begin{lstlisting}[style=pythonstyle]
x = [3, 1, 2]

print(sorted(x))
print(x)
\end{lstlisting}

Output:
\texttt{[1, 2, 3]}\\
\texttt{[3, 1, 2]}
\end{card}

% ============================================================
% PAGE 2
% ============================================================

\end{multicols}

\newpage

\begin{multicols}{2}
\raggedcolumns

% ------------------------------------------------------------
\section*{6. Containers}
% ------------------------------------------------------------

\begin{card}
\textbf{List} --- Ordered, mutable collection.

\begin{lstlisting}[style=pythonstyle]
x = [1, 2, 3]

x.append(4)
print(x)
\end{lstlisting}

Output:
\texttt{[1, 2, 3, 4]}
\end{card}

\begin{card}
\textbf{List methods}

\begin{lstlisting}[style=pythonstyle]
x = [1, 2, 3]

x.append(4)
x.extend([5, 6])
x.insert(0, 0)

print(x)
\end{lstlisting}

Output:
\texttt{[0, 1, 2, 3, 4, 5, 6]}
\end{card}

\begin{card}
\textbf{\texttt{pop()}} --- Remove and return an element.

\begin{lstlisting}[style=pythonstyle]
x = [10, 20, 30]

y = x.pop()

print(y)
print(x)
\end{lstlisting}

Output:
\texttt{30}\\
\texttt{[10, 20]}
\end{card}

\begin{card}
\textbf{\texttt{remove()}} --- Remove a specific value.

\begin{lstlisting}[style=pythonstyle]
x = [1, 2, 3, 2]

x.remove(2)

print(x)
\end{lstlisting}

Output:
\texttt{[1, 3, 2]}
\end{card}

\begin{card}
\textbf{\texttt{sort()}} --- Sort a list in place.

\begin{lstlisting}[style=pythonstyle]
x = [3, 1, 2]

x.sort()

print(x)
\end{lstlisting}

Output:
\texttt{[1, 2, 3]}
\end{card}

\begin{card}
\textbf{Tuple} --- Ordered, immutable collection.

\begin{lstlisting}[style=pythonstyle]
r = (1, 2, 3)

print(r)
print(r[0])
\end{lstlisting}

Output:
\texttt{(1, 2, 3)}\\
\texttt{1}

Useful for fixed quantities such as coordinates.
\end{card}

\begin{card}
\textbf{Unpacking} --- Assign several values at once.

\begin{lstlisting}[style=pythonstyle]
x = (10, 20, 30)

a, b, c = x

print(a)
print(b)
print(c)
\end{lstlisting}

Output:
\texttt{10}\\
\texttt{20}\\
\texttt{30}
\end{card}

\begin{card}
\textbf{Dictionary} --- Key--value storage.

\begin{lstlisting}[style=pythonstyle]
particle = {
    "mass": 2.0,
    "charge": 1.0
}

print(particle["mass"])
print(particle["charge"])
\end{lstlisting}

Output:
\texttt{2.0}\\
\texttt{1.0}
\end{card}

\begin{card}
\textbf{Dictionary methods}

\begin{lstlisting}[style=pythonstyle]
d = {"a": 1, "b": 2}

print(d.keys())
print(d.values())
print(d.items())
print(d.get("c", 0))
\end{lstlisting}

\texttt{get(key, default)} avoids an error if the key is absent.
\end{card}

\begin{card}
\textbf{Loop over a dictionary}

\begin{lstlisting}[style=pythonstyle]
d = {"m": 2, "q": 1}

for key, value in d.items():
    print(key, value)
\end{lstlisting}

Output:
\texttt{m 2}\\
\texttt{q 1}
\end{card}

\begin{card}
\textbf{Set} --- Unordered collection of unique values.

\begin{lstlisting}[style=pythonstyle]
x = {1, 2, 2, 3}

print(x)
print(2 in x)
\end{lstlisting}

Output:
\texttt{\{1, 2, 3\}}\\
\texttt{True}
\end{card}

% ------------------------------------------------------------
\section*{7. Indexing and Slicing}
% ------------------------------------------------------------

\begin{card}
\textbf{Indexing}

\begin{lstlisting}[style=pythonstyle]
x = [10, 20, 30, 40]

print(x[0])
print(x[-1])
\end{lstlisting}

Output:
\texttt{10}\\
\texttt{40}

Python starts indexing at zero.
\end{card}

\begin{card}
\textbf{Slicing}

\begin{lstlisting}[style=pythonstyle]
x = [0, 1, 2, 3, 4, 5]

print(x[1:4])
print(x[:3])
print(x[3:])
print(x[::2])
print(x[::-1])
\end{lstlisting}

Output:
\texttt{[1, 2, 3]}\\
\texttt{[0, 1, 2]}\\
\texttt{[3, 4, 5]}\\
\texttt{[0, 2, 4]}\\
\texttt{[5, 4, 3, 2, 1, 0]}
\end{card}

% ------------------------------------------------------------
\section*{8. \texttt{enumerate()} and \texttt{zip()}}
% ------------------------------------------------------------

\begin{card}
\textbf{\texttt{enumerate()}} --- Loop over index and value.

\begin{lstlisting}[style=pythonstyle]
x = [10, 20, 30]

for i, value in enumerate(x):
    print(i, value)
\end{lstlisting}

Output:
\texttt{0 10}\\
\texttt{1 20}\\
\texttt{2 30}
\end{card}

\begin{card}
\textbf{\texttt{zip()}} --- Iterate over sequences together.

\begin{lstlisting}[style=pythonstyle]
x = [1, 2, 3]
y = [4, 5, 6]

for a, b in zip(x, y):
    print(a, b)
\end{lstlisting}

Output:
\texttt{1 4}\\
\texttt{2 5}\\
\texttt{3 6}
\end{card}

% ------------------------------------------------------------
\section*{9. Comprehensions}
% ------------------------------------------------------------

\begin{card}
\textbf{List comprehension} --- Compact loop construction.

\begin{lstlisting}[style=pythonstyle]
x = [i**2 for i in range(5)]

print(x)
\end{lstlisting}

Output:
\texttt{[0, 1, 4, 9, 16]}
\end{card}

\begin{card}
\textbf{Conditional comprehension}

\begin{lstlisting}[style=pythonstyle]
x = [i for i in range(10) if i % 2 == 0]

print(x)
\end{lstlisting}

Output:
\texttt{[0, 2, 4, 6, 8]}
\end{card}

% ------------------------------------------------------------
\section*{10. Functions}
% ------------------------------------------------------------

\begin{card}
\textbf{\texttt{def}} --- Define a function.

\begin{lstlisting}[style=pythonstyle]
def energy(m, v):
    return 0.5 * m * v**2

E = energy(2, 3)

print(E)
\end{lstlisting}

Output: \texttt{9.0}
\end{card}

\begin{card}
\textbf{\texttt{return}} --- Send a value out of a function.

\begin{lstlisting}[style=pythonstyle]
def square(x):
    y = x**2
    return y

print(square(4))
\end{lstlisting}

Output: \texttt{16}
\end{card}

\begin{card}
\textbf{Multiple return values}

\begin{lstlisting}[style=pythonstyle]
def powers(x):
    return x**2, x**3

a, b = powers(2)

print(a)
print(b)
\end{lstlisting}

Output:
\texttt{4}\\
\texttt{8}
\end{card}

\begin{card}
\textbf{Default arguments}

\begin{lstlisting}[style=pythonstyle]
def greet(name, message="hello"):
    print(message, name)

greet("Chandra")
greet("Chandra", "hi")
\end{lstlisting}

Output:
\texttt{hello Chandra}\\
\texttt{hi Chandra}
\end{card}

\begin{card}
\textbf{\texttt{lambda}} --- Small anonymous function.

\begin{lstlisting}[style=pythonstyle]
f = lambda x: x**2

print(f(3))
\end{lstlisting}

Output: \texttt{9}

Often used by SciPy functions.
\end{card}

% ------------------------------------------------------------
\section*{11. References and Copying}
% ------------------------------------------------------------

\begin{card}
\textbf{Assignment does not copy a list.}

\begin{lstlisting}[style=pythonstyle]
a = [1, 2, 3]
b = a

b[0] = 100

print(a)
print(b)
\end{lstlisting}

Output:
\texttt{[100, 2, 3]}\\
\texttt{[100, 2, 3]}

Both names refer to the same list.
\end{card}

\begin{card}
\textbf{\texttt{copy()}} --- Make an independent copy.

\begin{lstlisting}[style=pythonstyle]
a = [1, 2, 3]
b = a.copy()

b[0] = 100

print(a)
print(b)
\end{lstlisting}

Output:
\texttt{[1, 2, 3]}\\
\texttt{[100, 2, 3]}
\end{card}

\begin{card}
\textbf{\texttt{is}} --- Test object identity.

\begin{lstlisting}[style=pythonstyle]
a = [1, 2]
b = a
c = a.copy()

print(a is b)
print(a is c)
\end{lstlisting}

Output:
\texttt{True}\\
\texttt{False}
\end{card}

% ============================================================
% PAGE 3
% ============================================================

\end{multicols}

\newpage

\begin{multicols}{2}
\raggedcolumns

% ------------------------------------------------------------
\section*{12. Modules and Imports}
% ------------------------------------------------------------

\begin{card}
\textbf{\texttt{import}} --- Load a module.

\begin{lstlisting}[style=pythonstyle]
import math

print(math.sqrt(16))
print(math.pi)
\end{lstlisting}

Output:
\texttt{4.0}\\
\texttt{3.141592653589793}
\end{card}

\begin{card}
\textbf{Import selected functions}

\begin{lstlisting}[style=pythonstyle]
from math import sqrt, pi

print(sqrt(16))
print(pi)
\end{lstlisting}
\end{card}

\begin{card}
\textbf{Alias a module}

\begin{lstlisting}[style=pythonstyle]
import numpy as np

x = np.array([1, 2, 3])

print(x)
\end{lstlisting}

Output:
\texttt{[1 2 3]}
\end{card}

% ------------------------------------------------------------
\section*{13. Useful Inspection Tools}
% ------------------------------------------------------------

\begin{card}
\textbf{\texttt{dir()}} --- See attributes and methods.

\begin{lstlisting}[style=pythonstyle]
x = [1, 2, 3]

print(dir(x))
\end{lstlisting}

Very useful when you forget a method.
\end{card}

\begin{card}
\textbf{\texttt{help()}} --- Read built-in documentation.

\begin{lstlisting}[style=pythonstyle]
help(str.split)
\end{lstlisting}

Useful at the interactive prompt.
\end{card}

\begin{card}
\textbf{\texttt{None}} --- Represents ``no value''.

\begin{lstlisting}[style=pythonstyle]
x = None

print(x)
print(x is None)
\end{lstlisting}

Output:
\texttt{None}\\
\texttt{True}
\end{card}

% ------------------------------------------------------------
\section*{14. Errors and Exceptions}
% ------------------------------------------------------------

\begin{card}
\textbf{\texttt{try / except}} --- Handle an error.

\begin{lstlisting}[style=pythonstyle]
try:
    x = 1 / 0
except ZeroDivisionError:
    print("division by zero")
\end{lstlisting}

Output:
\texttt{division by zero}
\end{card}

\begin{card}
\textbf{\texttt{raise}} --- Explicitly generate an error.

\begin{lstlisting}[style=pythonstyle]
def energy(m, v):
    if m <= 0:
        raise ValueError("mass must be positive")
    return 0.5 * m * v**2
\end{lstlisting}
\end{card}

\begin{card}
\textbf{\texttt{assert}} --- Check an assumption.

\begin{lstlisting}[style=pythonstyle]
mass = 2.0

assert mass > 0

print("valid")
\end{lstlisting}

Output:
\texttt{valid}

Useful for numerical sanity checks.
\end{card}

% ------------------------------------------------------------
\section*{15. Files}
% ------------------------------------------------------------

\begin{card}
\textbf{\texttt{open()}} --- Open a file.

\begin{lstlisting}[style=pythonstyle]
with open("data.txt", "r") as f:
    text = f.read()

print(text)
\end{lstlisting}

\texttt{"r"} = read.
\end{card}

\begin{card}
\textbf{Write a file}

\begin{lstlisting}[style=pythonstyle]
with open("output.txt", "w") as f:
    f.write("hello\n")
    f.write("world\n")
\end{lstlisting}

\texttt{"w"} = write and overwrite.
\end{card}

\begin{card}
\textbf{Read lines}

\begin{lstlisting}[style=pythonstyle]
with open("data.txt", "r") as f:
    lines = f.readlines()

print(lines)
\end{lstlisting}
\end{card}

\begin{card}
\textbf{\texttt{with}} --- Automatically closes the file.

\begin{lstlisting}[style=pythonstyle]
with open("data.txt") as f:
    data = f.read()
\end{lstlisting}

Prefer this to manually calling \texttt{close()}.
\end{card}

% ------------------------------------------------------------
\section*{16. NumPy Essentials}
% ------------------------------------------------------------

\begin{card}
\textbf{\texttt{np.array()}} --- Create numerical arrays.

\begin{lstlisting}[style=pythonstyle]
import numpy as np

x = np.array([1.0, 2.0, 3.0])

print(x)
print(type(x))
\end{lstlisting}

Output:
\texttt{[1. 2. 3.]}\\
\texttt{<class 'numpy.ndarray'>}
\end{card}

\begin{card}
\textbf{\texttt{np.zeros()}} --- Array of zeros.

\begin{lstlisting}[style=pythonstyle]
x = np.zeros(5)

print(x)
\end{lstlisting}

Output:
\texttt{[0. 0. 0. 0. 0.]}
\end{card}

\begin{card}
\textbf{\texttt{np.ones()}} --- Array of ones.

\begin{lstlisting}[style=pythonstyle]
x = np.ones(4)

print(x)
\end{lstlisting}

Output:
\texttt{[1. 1. 1. 1.]}
\end{card}

\begin{card}
\textbf{\texttt{np.eye()}} --- Identity matrix.

\begin{lstlisting}[style=pythonstyle]
I = np.eye(3)

print(I)
\end{lstlisting}
\end{card}

\begin{card}
\textbf{\texttt{np.linspace()}} --- Evenly spaced points.

\begin{lstlisting}[style=pythonstyle]
x = np.linspace(0, 1, 5)

print(x)
\end{lstlisting}

Output:
\texttt{[0. 0.25 0.5 0.75 1.]}

Very common for plotting.
\end{card}

\begin{card}
\textbf{\texttt{np.arange()}} --- Numerical range.

\begin{lstlisting}[style=pythonstyle]
x = np.arange(0, 1, 0.2)

print(x)
\end{lstlisting}

Output:
\texttt{[0. 0.2 0.4 0.6 0.8]}
\end{card}

\begin{card}
\textbf{\texttt{shape}} --- Array dimensions.

\begin{lstlisting}[style=pythonstyle]
A = np.zeros((3, 4))

print(A.shape)
print(A.ndim)
print(A.size)
\end{lstlisting}

Output:
\texttt{(3, 4)}\\
\texttt{2}\\
\texttt{12}
\end{card}

\begin{card}
\textbf{\texttt{reshape()}} --- Change array shape.

\begin{lstlisting}[style=pythonstyle]
x = np.arange(12)
A = x.reshape(3, 4)

print(A)
\end{lstlisting}
\end{card}

\begin{card}
\textbf{\texttt{dtype}} --- Numerical data type.

\begin{lstlisting}[style=pythonstyle]
x = np.array([1, 2, 3])

print(x.dtype)
\end{lstlisting}
\end{card}

% ============================================================
% PAGE 4
% ============================================================

\end{multicols}

\newpage

\begin{multicols}{2}
\raggedcolumns

% ------------------------------------------------------------
\section*{17. NumPy Indexing and Slicing}
% ------------------------------------------------------------

\begin{card}
\textbf{Array indexing}

\begin{lstlisting}[style=pythonstyle]
x = np.array([10, 20, 30, 40])

print(x[0])
print(x[-1])
\end{lstlisting}

Output:
\texttt{10}\\
\texttt{40}
\end{card}

\begin{card}
\textbf{Array slicing}

\begin{lstlisting}[style=pythonstyle]
x = np.arange(6)

print(x[1:4])
print(x[::2])
\end{lstlisting}

Output:
\texttt{[1 2 3]}\\
\texttt{[0 2 4]}
\end{card}

\begin{card}
\textbf{Matrix indexing}

\begin{lstlisting}[style=pythonstyle]
A = np.array([[1, 2],
              [3, 4]])

print(A[0, 1])
print(A[:, 0])
print(A[1, :])
\end{lstlisting}

Output:
\texttt{2}\\
\texttt{[1 3]}\\
\texttt{[3 4]}
\end{card}

% ------------------------------------------------------------
\section*{18. Vectorization and Broadcasting}
% ------------------------------------------------------------

\begin{card}
\textbf{Elementwise operations}

\begin{lstlisting}[style=pythonstyle]
x = np.array([1, 2, 3])

print(2 * x)
print(x**2)
print(x + 10)
\end{lstlisting}

Output:
\texttt{[2 4 6]}\\
\texttt{[1 4 9]}\\
\texttt{[11 12 13]}
\end{card}

\begin{card}
\textbf{Broadcasting} --- Compatible shapes can interact without explicit loops.

\begin{lstlisting}[style=pythonstyle]
x = np.array([1, 2, 3])

y = np.array([10, 20, 30])

print(x + y)
print(2 * x + y)
\end{lstlisting}

Output:
\texttt{[11 22 33]}\\
\texttt{[12 24 36]}
\end{card}

\begin{card}
\textbf{Boolean masking}

\begin{lstlisting}[style=pythonstyle]
x = np.array([-2, 1, 4, -1])

print(x > 0)
print(x[x > 0])
\end{lstlisting}

Output:
\texttt{[False True True False]}\\
\texttt{[1 4]}
\end{card}

\begin{card}
\textbf{\texttt{np.where()}} --- Conditional selection.

\begin{lstlisting}[style=pythonstyle]
x = np.array([-1, 2, -3, 4])

y = np.where(x > 0, x, 0)

print(y)
\end{lstlisting}

Output:
\texttt{[0 2 0 4]}
\end{card}

\begin{card}
\textbf{\texttt{np.all()} / \texttt{np.any()}}

\begin{lstlisting}[style=pythonstyle]
x = np.array([1, 2, 3])

print(np.all(x > 0))
print(np.any(x > 2))
\end{lstlisting}

Output:
\texttt{True}\\
\texttt{True}
\end{card}

% ------------------------------------------------------------
\section*{19. NumPy Numerical Tools}
% ------------------------------------------------------------

\begin{card}
\textbf{Elementary functions}

\begin{lstlisting}[style=pythonstyle]
x = np.array([0, 1, 2])

print(np.sqrt(x))
print(np.exp(x))
print(np.sin(x))
\end{lstlisting}
\end{card}

\begin{card}
\textbf{\texttt{np.sum()} / \texttt{np.mean()}}

\begin{lstlisting}[style=pythonstyle]
x = np.array([1, 2, 3, 4])

print(np.sum(x))
print(np.mean(x))
\end{lstlisting}

Output:
\texttt{10}\\
\texttt{2.5}
\end{card}

\begin{card}
\textbf{\texttt{np.std()} / \texttt{np.var()}}

\begin{lstlisting}[style=pythonstyle]
x = np.array([1, 2, 3, 4])

print(np.std(x))
print(np.var(x))
\end{lstlisting}
\end{card}

\begin{card}
\textbf{\texttt{np.argmax()} / \texttt{np.argmin()}} --- Index of maximum/minimum.

\begin{lstlisting}[style=pythonstyle]
x = np.array([4, 9, 2, 7])

print(np.argmax(x))
print(np.argmin(x))
\end{lstlisting}

Output:
\texttt{1}\\
\texttt{2}
\end{card}

\begin{card}
\textbf{\texttt{np.diff()}} --- Neighboring differences.

\begin{lstlisting}[style=pythonstyle]
x = np.array([1, 4, 9, 16])

print(np.diff(x))
\end{lstlisting}

Output:
\texttt{[3 5 7]}
\end{card}

\begin{card}
\textbf{\texttt{np.gradient()}} --- Numerical derivative.

\begin{lstlisting}[style=pythonstyle]
t = np.linspace(0, 10, 1000)
x = np.sin(t)

v = np.gradient(x, t)
\end{lstlisting}

Approximates
\[
v=\frac{dx}{dt}.
\]
\end{card}

\begin{card}
\textbf{\texttt{np.isclose()}} --- Compare floating-point values.

\begin{lstlisting}[style=pythonstyle]
print(np.isclose(0.1 + 0.2, 0.3))
\end{lstlisting}

Output: \texttt{True}

Prefer this to \texttt{==} for floating-point comparisons.
\end{card}

\begin{card}
\textbf{\texttt{np.allclose()}} --- Compare arrays.

\begin{lstlisting}[style=pythonstyle]
A = np.array([1.0, 2.0])
B = np.array([1.0, 2.0000001])

print(np.allclose(A, B))
\end{lstlisting}

Output: \texttt{True}
\end{card}

\begin{card}
\textbf{\texttt{np.isnan()} / \texttt{np.isfinite()}}

\begin{lstlisting}[style=pythonstyle]
x = np.array([1.0, np.nan, np.inf])

print(np.isnan(x))
print(np.isfinite(x))
\end{lstlisting}
\end{card}

% ------------------------------------------------------------
\section*{20. Linear Algebra}
% ------------------------------------------------------------

\begin{card}
\textbf{\texttt{np.dot()}} --- Dot product.

\begin{lstlisting}[style=pythonstyle]
a = np.array([1, 2, 3])
b = np.array([4, 5, 6])

print(np.dot(a, b))
\end{lstlisting}

Output:
\texttt{32}

\[
\mathbf a\cdot\mathbf b.
\]
\end{card}

\begin{card}
\textbf{\texttt{@}} --- Matrix multiplication.

\begin{lstlisting}[style=pythonstyle]
A = np.array([[1, 2],
              [3, 4]])

B = np.array([[5, 6],
              [7, 8]])

print(A @ B)
\end{lstlisting}

Output:
\[
\begin{pmatrix}
19 & 22\\
43 & 50
\end{pmatrix}.
\]
\end{card}

\begin{card}
\textbf{\texttt{.T}} --- Transpose.

\begin{lstlisting}[style=pythonstyle]
A = np.array([[1, 2],
              [3, 4]])

print(A.T)
\end{lstlisting}
\end{card}

\begin{card}
\textbf{\texttt{np.cross()}} --- Cross product.

\begin{lstlisting}[style=pythonstyle]
r = np.array([1, 0, 0])
p = np.array([0, 2, 0])

L = np.cross(r, p)

print(L)
\end{lstlisting}

Output:
\texttt{[0 0 2]}

\[
\mathbf L=\mathbf r\times\mathbf p.
\]
\end{card}

\begin{card}
\textbf{\texttt{np.linalg.norm()}} --- Vector magnitude.

\begin{lstlisting}[style=pythonstyle]
r = np.array([3, 4])

print(np.linalg.norm(r))
\end{lstlisting}

Output: \texttt{5.0}
\end{card}

\begin{card}
\textbf{Unit vector}

\begin{lstlisting}[style=pythonstyle]
r_hat = r / np.linalg.norm(r)

print(r_hat)
\end{lstlisting}

\[
\hat{\mathbf r}
=
\frac{\mathbf r}{|\mathbf r|}.
\]
\end{card}

% ============================================================
% PAGE 5
% ============================================================

\end{multicols}

\newpage

\begin{multicols}{2}
\raggedcolumns

% ------------------------------------------------------------
\section*{21. More Linear Algebra}
% ------------------------------------------------------------

\begin{card}
\textbf{\texttt{np.linalg.solve()}} --- Solve
\[
A\mathbf x=\mathbf b.
\]

\begin{lstlisting}[style=pythonstyle]
A = np.array([[2, 1],
              [1, 3]])

b = np.array([5, 6])

x = np.linalg.solve(A, b)

print(x)
\end{lstlisting}
\end{card}

\begin{card}
\textbf{\texttt{np.linalg.eig()}} --- Eigenvalues/eigenvectors.

\begin{lstlisting}[style=pythonstyle]
A = np.array([[2, 0],
              [0, 3]])

w, V = np.linalg.eig(A)

print(w)
print(V)
\end{lstlisting}

\[
A\mathbf v=\lambda\mathbf v.
\]
\end{card}

\begin{card}
\textbf{\texttt{np.linalg.det()}} --- Determinant.

\begin{lstlisting}[style=pythonstyle]
A = np.array([[1, 2],
              [3, 4]])

print(np.linalg.det(A))
\end{lstlisting}

Output:
\texttt{-2.0}
\end{card}

% ------------------------------------------------------------
\section*{22. Matplotlib}
% ------------------------------------------------------------

\begin{card}
\textbf{\texttt{plt.plot()}} --- Line plot.

\begin{lstlisting}[style=pythonstyle]
import matplotlib.pyplot as plt

t = np.linspace(0, 10, 1000)
x = np.cos(t)

plt.plot(t, x)
plt.show()
\end{lstlisting}
\end{card}

\begin{card}
\textbf{Labels and title}

\begin{lstlisting}[style=pythonstyle]
plt.xlabel("Time")
plt.ylabel("Position")
plt.title("Oscillator")
\end{lstlisting}
\end{card}

\begin{card}
\textbf{\texttt{plt.legend()}}

\begin{lstlisting}[style=pythonstyle]
plt.plot(t, np.cos(t), label="cos")
plt.plot(t, np.sin(t), label="sin")

plt.legend()
plt.show()
\end{lstlisting}
\end{card}

\begin{card}
\textbf{\texttt{plt.grid()}}

\begin{lstlisting}[style=pythonstyle]
plt.plot(t, x)

plt.grid()
plt.show()
\end{lstlisting}
\end{card}

\begin{card}
\textbf{\texttt{plt.scatter()}} --- Scatter plot.

\begin{lstlisting}[style=pythonstyle]
plt.scatter(t, x)
plt.show()
\end{lstlisting}

Useful for data points.
\end{card}

\begin{card}
\textbf{\texttt{plt.subplots()}} --- Explicit figure and axes.

\begin{lstlisting}[style=pythonstyle]
fig, ax = plt.subplots()

ax.plot(t, x)
ax.set_xlabel("t")
ax.set_ylabel("x")

plt.show()
\end{lstlisting}

Preferred for larger programs.
\end{card}

\begin{card}
\textbf{Logarithmic axes}

\begin{lstlisting}[style=pythonstyle]
plt.loglog(x, y)
plt.show()
\end{lstlisting}

Also:
\texttt{plt.semilogx()} and
\texttt{plt.semilogy()}.
\end{card}

\begin{card}
\textbf{\texttt{plt.savefig()}} --- Save a figure.

\begin{lstlisting}[style=pythonstyle]
plt.savefig("spectrum.png", dpi=300)
\end{lstlisting}
\end{card}

% ------------------------------------------------------------
\section*{23. SciPy}
% ------------------------------------------------------------

\begin{card}
\textbf{\texttt{solve\_ivp()}} --- Solve an ODE.

\begin{lstlisting}[style=pythonstyle]
from scipy.integrate import solve_ivp

def rhs(t, y):
    x, v = y
    return [v, -x]

sol = solve_ivp(
    rhs,
    [0, 10],
    [1, 0]
)

print(sol.t.shape)
print(sol.y.shape)
\end{lstlisting}

The state is
\[
\mathbf y=
\begin{pmatrix}
x\\
v
\end{pmatrix}.
\]
\end{card}

\begin{card}
\textbf{\texttt{quad()}} --- Numerical integration.

\begin{lstlisting}[style=pythonstyle]
from scipy.integrate import quad

f = lambda x: x**2

I, err = quad(f, 0, 1)

print(I)
print(err)
\end{lstlisting}
\end{card}

\begin{card}
\textbf{\texttt{root\_scalar()}} --- Find a scalar root.

\begin{lstlisting}[style=pythonstyle]
from scipy.optimize import root_scalar

f = lambda x: x**2 - 2

result = root_scalar(
    f,
    bracket=[0, 2]
)

print(result.root)
\end{lstlisting}

Output is approximately:
\texttt{1.41421356}
\end{card}

\begin{card}
\textbf{\texttt{minimize()}} --- Numerical optimization.

\begin{lstlisting}[style=pythonstyle]
from scipy.optimize import minimize

f = lambda x: (x - 3)**2

result = minimize(f, 0)

print(result.x)
\end{lstlisting}

Output is approximately:
\texttt{[3.]}
\end{card}

\begin{card}
\textbf{\texttt{curve\_fit()}} --- Fit a model to data.

\begin{lstlisting}[style=pythonstyle]
from scipy.optimize import curve_fit

def model(x, a, b):
    return a*x + b

popt, pcov = curve_fit(
    model, x, y
)

print(popt)
\end{lstlisting}

\texttt{popt} contains fitted parameters.
\texttt{pcov} contains their covariance estimate.
\end{card}

% ------------------------------------------------------------
\section*{24. SymPy}
% ------------------------------------------------------------

\begin{card}
\textbf{\texttt{symbols()}} --- Create symbolic variables.

\begin{lstlisting}[style=pythonstyle]
import sympy as sp

x = sp.symbols("x")

print(x)
\end{lstlisting}

Output: \texttt{x}
\end{card}

\begin{card}
\textbf{\texttt{diff()}} --- Symbolic differentiation.

\begin{lstlisting}[style=pythonstyle]
f = x**3 + 2*x

print(sp.diff(f, x))
\end{lstlisting}

Output:
\[
3x^2+2.
\]
\end{card}

\begin{card}
\textbf{\texttt{integrate()}} --- Symbolic integration.

\begin{lstlisting}[style=pythonstyle]
print(
    sp.integrate(x**2, x)
)
\end{lstlisting}

Output:
\[
\frac{x^3}{3}.
\]
\end{card}

\begin{card}
\textbf{\texttt{solve()}} --- Solve equations symbolically.

\begin{lstlisting}[style=pythonstyle]
print(
    sp.solve(x**2 - 4, x)
)
\end{lstlisting}

Output:
\texttt{[-2, 2]}
\end{card}

\begin{card}
\textbf{\texttt{simplify()}} --- Simplify expressions.

\begin{lstlisting}[style=pythonstyle]
expr = (x**2 - 1)/(x - 1)

print(sp.simplify(expr))
\end{lstlisting}

Result:
\[
x+1.
\]
\end{card}

\begin{card}
\textbf{\texttt{series()}} --- Taylor expansion.

\begin{lstlisting}[style=pythonstyle]
print(
    sp.series(sp.sin(x), x, 0, 5)
)
\end{lstlisting}

Output:
\[
x-\frac{x^3}{6}+O(x^5).
\]
\end{card}

% ============================================================
% PAGE 6
% ============================================================

\end{multicols}

\newpage

\begin{multicols}{2}
\raggedcolumns

% ------------------------------------------------------------
\section*{25. Useful Scientific Patterns}
% ------------------------------------------------------------

\begin{card}
\textbf{Typical physics workflow}

\begin{lstlisting}[style=pythonstyle]
import numpy as np
import matplotlib.pyplot as plt

t = np.linspace(0, 10, 1000)
x = np.cos(t)

plt.plot(t, x)
plt.xlabel("t")
plt.ylabel("x(t)")
plt.show()
\end{lstlisting}
\end{card}

\begin{card}
\textbf{Parameter sweep}

\begin{lstlisting}[style=pythonstyle]
for alpha in [0.1, 0.2, 0.5]:
    result = simulate(alpha)
    print(alpha, result)
\end{lstlisting}

Useful for scanning model parameters.
\end{card}

\begin{card}
\textbf{Numerical derivative}

\begin{lstlisting}[style=pythonstyle]
v = np.gradient(x, t)
a = np.gradient(v, t)

print(v)
print(a)
\end{lstlisting}

Corresponds to
\[
v=\dot{x},
\qquad
a=\ddot{x}.
\]
\end{card}

\begin{card}
\textbf{Conservation check}

\begin{lstlisting}[style=pythonstyle]
E = kinetic + potential

print(np.max(E))
print(np.min(E))
print(
    np.max(E) - np.min(E)
)
\end{lstlisting}

The last quantity measures the numerical variation of $E$.
\end{card}

\begin{card}
\textbf{Position, velocity, acceleration}

\begin{lstlisting}[style=pythonstyle]
r = np.array([x, y, z])
v = np.array([vx, vy, vz])

speed = np.linalg.norm(v)

print(speed)
\end{lstlisting}

\[
|\mathbf v|
=
\sqrt{\mathbf v\cdot\mathbf v}.
\]
\end{card}

\begin{card}
\textbf{Angular momentum}

\begin{lstlisting}[style=pythonstyle]
L = np.cross(r, p)

print(L)
\end{lstlisting}

\[
\mathbf L=\mathbf r\times\mathbf p.
\]
\end{card}

\begin{card}
\textbf{Energy}

\begin{lstlisting}[style=pythonstyle]
K = 0.5 * m * np.dot(v, v)
U = potential(r)

E = K + U

print(K)
print(U)
print(E)
\end{lstlisting}

\[
E=K+U.
\]
\end{card}

\begin{card}
\textbf{ODE state vector}

For
\[
\ddot{x}=f(x,\dot{x},t),
\]
use
\[
\mathbf y=
\begin{pmatrix}
x\\
\dot{x}
\end{pmatrix}.
\]

\begin{lstlisting}[style=pythonstyle]
def rhs(t, y):
    x, v = y
    return [v, force(x, v, t)]
\end{lstlisting}
\end{card}

% ------------------------------------------------------------
\section*{26. Random Numbers}
% ------------------------------------------------------------

\begin{card}
\textbf{Modern random-number generator}

\begin{lstlisting}[style=pythonstyle]
rng = np.random.default_rng(123)

x = rng.random(5)

print(x)
\end{lstlisting}

A seed makes the sequence reproducible.
\end{card}

\begin{card}
\textbf{Gaussian random numbers}

\begin{lstlisting}[style=pythonstyle]
rng = np.random.default_rng(123)

x = rng.normal(
    loc=0,
    scale=1,
    size=5
)

print(x)
\end{lstlisting}

Useful for simulations and Monte Carlo methods.
\end{card}

% ------------------------------------------------------------
\section*{27. Script Structure}
% ------------------------------------------------------------

\begin{card}
\textbf{\texttt{if \_\_name\_\_ == "\_\_main\_\_"}}

\begin{lstlisting}[style=pythonstyle]
def main():
    print("running program")

if __name__ == "__main__":
    main()
\end{lstlisting}

The block runs when the file is executed directly,
not when it is imported as a module.
\end{card}

\begin{card}
\textbf{Typical reusable script}

\begin{lstlisting}[style=pythonstyle]
import numpy as np

def model(x):
    return np.exp(-x**2)

def main():
    x = np.linspace(-5, 5, 1000)
    y = model(x)
    print(y[0])

if __name__ == "__main__":
    main()
\end{lstlisting}
\end{card}

% ------------------------------------------------------------
\section*{28. What to Memorize First}
% ------------------------------------------------------------

\begin{card}
\begin{itemize}[leftmargin=*,nosep]
    \item Variables, numbers, strings, booleans, \texttt{None}
    \item \texttt{+ - * / // \% **}
    \item \texttt{== != < > <= >=}
    \item \texttt{and}, \texttt{or}, \texttt{not}, \texttt{in}
    \item \texttt{if / elif / else}
    \item \texttt{for}, \texttt{while}
    \item \texttt{break}, \texttt{continue}
    \item Lists, tuples, dictionaries, sets
    \item Indexing and slicing
    \item \texttt{enumerate()}, \texttt{zip()}
    \item Functions, arguments, \texttt{return}
    \item List comprehensions
    \item References vs.\ copies
    \item \texttt{import}
    \item \texttt{try / except}, \texttt{assert}
    \item File handling with \texttt{with open()}
    \item NumPy arrays and \texttt{shape}
    \item Vectorization and broadcasting
    \item \texttt{dot()}, \texttt{cross()}, \texttt{norm()}
    \item Matplotlib plotting
    \item SciPy ODE solving, integration, optimization
    \item SymPy differentiation and solving
\end{itemize}
\end{card}

\end{multicols}

% ============================================================
% PAGE 7+ : the rest of what "Statistics and Data Analysis for the
% Physicist" uses (chosen from a census of the book's 343 scripts)
% ============================================================

\newpage

\begin{center}
    {\Large\bfseries\color{navy}
    The Rest of What the Book Uses}\\[2pt]
    {\small Sections 29--41: everything the book's scripts call that sections 1--28 do not cover}
\end{center}

\begin{multicols}{2}
\raggedcolumns

% ------------------------------------------------------------
\section*{29. The Top of Every Script}
% ------------------------------------------------------------

\begin{card}
\textbf{\texttt{pathlib.Path}} --- Paths that work on every computer.

\begin{lstlisting}[style=pythonstyle]
from pathlib import Path

here = Path(__file__).resolve()
print(here.name)        # this file
print(here.parent)      # its folder
print(here.parents[1])  # one folder up
\end{lstlisting}

\texttt{resolve()} turns a relative path into a full one.
\end{card}

\begin{card}
\textbf{\texttt{sys.path}} --- Let Python find the book's shared code.

\begin{lstlisting}[style=pythonstyle]
import sys
from pathlib import Path

sys.path.insert(0, str(Path(__file__).resolve().parents[1]))
from common import setup, savefig, save_numbers
\end{lstlisting}

Adds \texttt{code/} to the import search, so \texttt{common.py} can be
imported from any chapter folder.
\end{card}

\begin{card}
\textbf{Files on disk}

\begin{lstlisting}[style=pythonstyle]
out = Path("data/ch06")
out.mkdir(parents=True, exist_ok=True)

f = out / "samples.npz"
print(f.exists())
\end{lstlisting}

The \texttt{/} operator joins path pieces.
\end{card}

% ------------------------------------------------------------
\section*{30. Converting Types}
% ------------------------------------------------------------

\begin{card}
\textbf{\texttt{int()}, \texttt{float()}, \texttt{str()}}

\begin{lstlisting}[style=pythonstyle]
print(int(3.9))
print(int("42"))
print(float("1e-3"))
print(str(3.0) + " m")
\end{lstlisting}

Output:
\texttt{3}\\
\texttt{42}\\
\texttt{0.001}\\
\texttt{3.0 m}

\texttt{int()} cuts towards zero; it does not round.
\end{card}

\begin{card}
\textbf{\texttt{round()}}

\begin{lstlisting}[style=pythonstyle]
print(round(3.14159, 2))
print(round(2.675, 2))
\end{lstlisting}

Output:
\texttt{3.14}\\
\texttt{2.67}

$2.675$ is stored as $2.67499\ldots$ in binary, so it rounds down.
\end{card}

\begin{card}
\textbf{\texttt{.astype()}} --- Convert a whole array.

\begin{lstlisting}[style=pythonstyle]
x = np.array([1.7, -1.7])

print(x.astype(int))
\end{lstlisting}

Output:
\texttt{[ 1 -1]}
\end{card}

\begin{card}
\textbf{\texttt{dict.update()}} --- Add or overwrite several keys.

\begin{lstlisting}[style=pythonstyle]
d = {"a": 1, "b": 2}
d.update({"b": 3, "c": 4})

print(d)
\end{lstlisting}

Output:
\texttt{\{'a': 1, 'b': 3, 'c': 4\}}
\end{card}

% ------------------------------------------------------------
\section*{31. Matplotlib the Book's Way: \texttt{fig, ax}}
% ------------------------------------------------------------

\begin{card}
\textbf{One figure, one axes} --- every call goes to \texttt{ax}.

\begin{lstlisting}[style=pythonstyle]
fig, ax = plt.subplots(figsize=(5, 3.5))

ax.plot(x, y, label="model")
ax.set_title("Power spectrum")
ax.set_xlim(0, 10)
ax.set_ylim(-1, 1)
ax.legend()

fig.tight_layout()
fig.savefig("spectrum.pdf")
\end{lstlisting}

\texttt{tight\_layout()} stops labels from overlapping.
\end{card}

\begin{card}
\textbf{Several panels}

\begin{lstlisting}[style=pythonstyle]
fig, axes = plt.subplots(1, 2, figsize=(8, 3))

axes[0].plot(x, y)
axes[1].hist(samples, bins=40)
\end{lstlisting}

\texttt{axes} is an array of axes, one per panel.
\end{card}

\begin{card}
\textbf{Reference lines and text}

\begin{lstlisting}[style=pythonstyle]
ax.axhline(0, color="gray", lw=0.8)
ax.axvline(1.96, ls="--")
ax.text(0.05, 0.9, "p = 0.05",
        transform=ax.transAxes)
\end{lstlisting}

With \texttt{transform=ax.transAxes} the position is a fraction of the
panel, $(0,0)$ bottom left.
\end{card}

\begin{card}
\textbf{Histograms and error bars}

\begin{lstlisting}[style=pythonstyle]
ax.hist(samples, bins=50, density=True)

ax.errorbar(x, y, yerr=sigma, fmt="o",
            capsize=2)
\end{lstlisting}

\texttt{density=True} makes the bars integrate to one, comparable with a pdf.
\end{card}

\begin{card}
\textbf{Bands, scales, images}

\begin{lstlisting}[style=pythonstyle]
ax.fill_between(x, lo, hi, alpha=0.3)

ax.set_xscale("log")
ax.set_yscale("log")

im = ax.imshow(Z, origin="lower")
fig.colorbar(im, ax=ax)
ax.set_aspect("equal")
\end{lstlisting}

\texttt{fill\_between} draws error bands such as $\pm1\sigma$.
\end{card}

% ------------------------------------------------------------
\section*{32. Building Arrays}
% ------------------------------------------------------------

\begin{card}
\textbf{\texttt{np.full()}, \texttt{np.empty()}, \texttt{np.zeros\_like()}}

\begin{lstlisting}[style=pythonstyle]
print(np.full(3, 7.0))

a = np.empty(1000)   # filled later
b = np.zeros_like(a) # same shape, zeros
\end{lstlisting}

Output:
\texttt{[7. 7. 7.]}

\texttt{np.empty} contains leftover garbage until you fill it.
\end{card}

\begin{card}
\textbf{\texttt{np.logspace()} / \texttt{np.geomspace()}} --- Evenly spaced in log.

\begin{lstlisting}[style=pythonstyle]
print(np.geomspace(1, 1000, 4))
print(np.logspace(0, 3, 4))
\end{lstlisting}

Both give $1, 10, 100, 1000$: \texttt{geomspace} takes the end values,
\texttt{logspace} their powers of ten.
\end{card}

\begin{card}
\textbf{\texttt{np.diag()}}

\begin{lstlisting}[style=pythonstyle]
C = np.diag([1.0, 4.0])  # matrix from diagonal
print(np.diag(C))        # diagonal from matrix
\end{lstlisting}

Output:
\texttt{[1. 4.]}

Error bars from a covariance: \texttt{np.sqrt(np.diag(C))}.
\end{card}

\begin{card}
\textbf{Joining arrays}

\begin{lstlisting}[style=pythonstyle]
a = np.array([1, 2])
b = np.array([3, 4])

print(np.concatenate([a, b]))
print(np.column_stack([a, b]))
\end{lstlisting}

Output:
\texttt{[1 2 3 4]}\\
\texttt{[[1 3]}\\
\texttt{\ [2 4]]}

Also \texttt{np.vstack} (rows) and \texttt{np.stack} (new axis).
\end{card}

\begin{card}
\textbf{\texttt{.ravel()}} --- Flatten to one dimension.

\begin{lstlisting}[style=pythonstyle]
A = np.array([[1, 2], [3, 4]])

print(A.ravel())
\end{lstlisting}

Output:
\texttt{[1 2 3 4]}
\end{card}

\begin{card}
\textbf{\texttt{np.meshgrid()}} --- A 2D grid of coordinates.

\begin{lstlisting}[style=pythonstyle]
x = np.linspace(-1, 1, 200)
y = np.linspace(-1, 1, 100)
X, Y = np.meshgrid(x, y)

R = np.sqrt(X**2 + Y**2)
print(R.shape)
\end{lstlisting}

Output:
\texttt{(100, 200)}

Rows follow $y$, columns follow $x$.
\end{card}

% ------------------------------------------------------------
\section*{33. More NumPy Maths}
% ------------------------------------------------------------

\begin{card}
\textbf{Logs, powers, absolute values}

\begin{lstlisting}[style=pythonstyle]
print(np.log(np.e))
print(np.log10(1000))
print(np.abs(3 + 4j))
\end{lstlisting}

Output:
\texttt{1.0}\\
\texttt{3.0}\\
\texttt{5.0}

\texttt{np.log} is the natural log, $\ln$.
\end{card}

\begin{card}
\textbf{\texttt{np.cumsum()}, \texttt{np.maximum()}, \texttt{np.clip()}}

\begin{lstlisting}[style=pythonstyle]
print(np.cumsum([1, 2, 3]))
print(np.maximum([1, 5], [3, 2]))
print(np.clip([-2, 0.5, 3], 0, 1))
\end{lstlisting}

Output:
\texttt{[1 3 6]}\\
\texttt{[3 5]}\\
\texttt{[0.\ \ 0.5 1.\ ]}

\texttt{np.maximum} compares element by element; \texttt{np.max} gives one number.
\end{card}

\begin{card}
\textbf{\texttt{np.median()}, \texttt{np.percentile()}, \texttt{np.quantile()}}

\begin{lstlisting}[style=pythonstyle]
x = rng.standard_normal(100000)

print(np.median(x))
print(np.percentile(x, [16, 84]))
print(np.quantile(x, 0.975))
\end{lstlisting}

Close to $0$, $[-1, 1]$ and $1.96$: the $68\%$ interval and a one-sided
$97.5\%$ point of a Gaussian.
\end{card}

\begin{card}
\textbf{\texttt{np.interp()}} --- Straight-line interpolation.

\begin{lstlisting}[style=pythonstyle]
print(np.interp(2.5, [1, 2, 3], [10, 20, 30]))
\end{lstlisting}

Output: \texttt{25.0}

\texttt{x} must be increasing.
\end{card}

\begin{card}
\textbf{\texttt{np.trapezoid()}} --- Integrate sampled values.

\begin{lstlisting}[style=pythonstyle]
x = np.linspace(0, 1, 101)

print(np.trapezoid(x**2, x))
\end{lstlisting}

Output: about \texttt{0.33335}, against the exact $1/3$.

Called \texttt{np.trapz} before NumPy 2.0. Use \texttt{quad} when you
have the function, \texttt{trapezoid} when you only have samples.
\end{card}

\begin{card}
\textbf{Counting: \texttt{np.histogram()}, \texttt{np.bincount()}}

\begin{lstlisting}[style=pythonstyle]
counts, edges = np.histogram(x, bins=20)
print(len(counts), len(edges))

print(np.bincount([0, 1, 1, 3]))
\end{lstlisting}

Output:
\texttt{20 21}\\
\texttt{[1 2 0 1]}

$N$ bins have $N+1$ edges.
\end{card}

\begin{card}
\textbf{Sorting and searching}

\begin{lstlisting}[style=pythonstyle]
x = np.array([30, 10, 20])

print(np.sort(x))
print(np.argsort(x))
print(np.searchsorted([1, 3, 5], 4))
print(np.unique([3, 1, 3, 2]))
\end{lstlisting}

Output:
\texttt{[10 20 30]}\\
\texttt{[1 2 0]}\\
\texttt{2}\\
\texttt{[1 2 3]}
\end{card}

\begin{card}
\textbf{\texttt{np.cov()}} --- Sample covariance matrix.

\begin{lstlisting}[style=pythonstyle]
C = np.cov(x, y)   # 2 x 2
print(C[0, 1])     # covariance of x and y
\end{lstlisting}

Divides by $N-1$. Rows are variables.
\end{card}

\begin{card}
\textbf{\texttt{np.outer()} and \texttt{np.einsum()}}

\begin{lstlisting}[style=pythonstyle]
print(np.outer([1, 2], [3, 4]))

chi2 = np.einsum("i,ij,j->", r, Cinv, r)
\end{lstlisting}

Output:
\texttt{[[3 4]}\\
\texttt{\ [6 8]]}

The \texttt{einsum} line is $\chi^2=\mathbf r^{\mathsf T}C^{-1}\mathbf r$:
repeated letters are summed.
\end{card}

\begin{card}
\textbf{\texttt{np.polyfit()}} --- Least-squares polynomial.

\begin{lstlisting}[style=pythonstyle]
slope, intercept = np.polyfit(x, y, 1)
\end{lstlisting}

Highest power first.
\end{card}

% ------------------------------------------------------------
\section*{34. More Random Numbers}
% ------------------------------------------------------------

\begin{card}
\textbf{The generator's other methods}

\begin{lstlisting}[style=pythonstyle]
rng = np.random.default_rng(42)

z = rng.standard_normal(1000)    # N(0, 1)
u = rng.uniform(0, 1, size=1000)
n = rng.poisson(4.0, size=1000)  # counts
k = rng.integers(0, 10, size=5)  # 0..9
\end{lstlisting}

\texttt{integers(0, 10)} never returns $10$.
\end{card}

\begin{card}
\textbf{\texttt{choice()}} --- Draw from a list (resampling).

\begin{lstlisting}[style=pythonstyle]
boot = rng.choice(data, size=len(data),
                  replace=True)
\end{lstlisting}

One bootstrap sample. \texttt{replace=False} shuffles without repeats.
\end{card}

\begin{card}
\textbf{\texttt{multivariate\_normal()}} --- Correlated Gaussians.

\begin{lstlisting}[style=pythonstyle]
mean = [0, 0]
C = [[1.0, 0.8],
     [0.8, 1.0]]

xy = rng.multivariate_normal(mean, C, size=5000)
print(xy.shape)
\end{lstlisting}

Output:
\texttt{(5000, 2)}
\end{card}

% ------------------------------------------------------------
\section*{35. Saving and Loading Arrays}
% ------------------------------------------------------------

\begin{card}
\textbf{\texttt{np.savez()} / \texttt{np.load()}}

\begin{lstlisting}[style=pythonstyle]
np.savez("run.npz", ell=ell, cl=cl)

d = np.load("run.npz")
print(d["cl"][:3])
\end{lstlisting}

Several named arrays in one file, read back exactly.
\end{card}

\begin{card}
\textbf{\texttt{np.loadtxt()}} --- Read a table of numbers.

\begin{lstlisting}[style=pythonstyle]
z, H, sigma = np.loadtxt("chronometers.txt",
                         unpack=True)
\end{lstlisting}

\texttt{unpack=True} gives one array per column. Lines starting with
\texttt{\#} are skipped.
\end{card}

% ------------------------------------------------------------
\section*{36. Linear Algebra for Statistics}
% ------------------------------------------------------------

\begin{card}
\textbf{\texttt{solve}, not \texttt{inv}}

\begin{lstlisting}[style=pythonstyle]
x = np.linalg.solve(C, b)    # good
x = np.linalg.inv(C) @ b     # same, less accurate
\end{lstlisting}

Use \texttt{inv} only when the inverse itself is wanted, e.g.\ the error
matrix $F^{-1}$ from a Fisher matrix $F$.
\end{card}

\begin{card}
\textbf{\texttt{np.linalg.cholesky()}} --- $C=LL^{\mathsf T}$.

\begin{lstlisting}[style=pythonstyle]
L = np.linalg.cholesky(C)

z = rng.standard_normal(len(C))
x = L @ z            # covariance C

logdet = 2 * np.sum(np.log(np.diag(L)))
\end{lstlisting}

Correlated noise and $\ln\det C$ in one step. Fails if $C$ is not
positive definite: a useful warning.
\end{card}

\begin{card}
\textbf{\texttt{np.linalg.eigh()}} --- Symmetric matrices.

\begin{lstlisting}[style=pythonstyle]
w, V = np.linalg.eigh(C)

print(w)          # ascending
print(w.max() / w.min())
\end{lstlisting}

The last line is the condition number: about how many times errors in
$\mathbf b$ can be magnified when solving $C\mathbf x=\mathbf b$.
\end{card}

% ------------------------------------------------------------
\section*{37. Fourier Transforms}
% ------------------------------------------------------------

\begin{card}
\textbf{\texttt{np.fft.rfft()}} --- Spectrum of a real signal.

\begin{lstlisting}[style=pythonstyle]
n, dt = 1024, 0.01
t = np.arange(n) * dt
x = np.sin(2 * np.pi * 5 * t)

X = np.fft.rfft(x)
f = np.fft.rfftfreq(n, dt)

print(f[np.argmax(np.abs(X))])
\end{lstlisting}

Output: about \texttt{4.98}, the nearest frequency bin to the $5$\,Hz input.

Power: \texttt{np.abs(X)**2}. Back to time: \texttt{np.fft.irfft(X, n)}.
\end{card}

% ------------------------------------------------------------
\section*{38. Probability Distributions: \texttt{scipy.stats}}
% ------------------------------------------------------------

\begin{card}
\textbf{The four questions to ask a distribution}

\begin{lstlisting}[style=pythonstyle]
from scipy import stats

print(stats.norm.pdf(0))      # density
print(stats.norm.cdf(1.96))   # P(X <= x)
print(stats.norm.sf(5))       # P(X > x)
print(stats.norm.ppf(0.975))  # inverse cdf
\end{lstlisting}

Output:
\texttt{0.3989...}\\
\texttt{0.9750...}\\
\texttt{2.8665...e-07}\\
\texttt{1.9599...}

\texttt{sf} $=1-$\texttt{cdf}, but accurate far in the tail (the $5\sigma$
$p$-value).
\end{card}

\begin{card}
\textbf{Other distributions, same methods}

\begin{lstlisting}[style=pythonstyle]
print(stats.chi2.sf(30, df=20))   # p-value
print(stats.chi2.ppf(0.95, df=1))
print(stats.poisson.pmf(3, 2.0))  # P(N = 3)

x = stats.norm.rvs(loc=1, scale=2,
                   size=5, random_state=rng)
\end{lstlisting}

Output (first three):
about \texttt{0.07}, \texttt{3.84}, \texttt{0.180}.

Discrete distributions have \texttt{pmf} instead of \texttt{pdf}.
\end{card}

% ------------------------------------------------------------
\section*{39. Roots and Special Functions}
% ------------------------------------------------------------

\begin{card}
\textbf{\texttt{brentq()}} --- The book's root finder.

\begin{lstlisting}[style=pythonstyle]
from scipy.optimize import brentq

r = brentq(lambda x: x**2 - 2, 0, 2)

print(r)
\end{lstlisting}

Output: \texttt{1.414213562373...}

$f$ must change sign between the two ends.
\end{card}

\begin{card}
\textbf{\texttt{gammaln()} / \texttt{logsumexp()}} --- Work in logs.

\begin{lstlisting}[style=pythonstyle]
from scipy.special import gammaln, logsumexp

print(gammaln(101))           # ln(100!)
print(logsumexp([1000, 1000]))
\end{lstlisting}

Output:
\texttt{363.739...}\\
\texttt{1000.693...}

$100!$ and $e^{1000}$ overflow; their logarithms do not.
\end{card}

% ------------------------------------------------------------
\section*{40. Timing}
% ------------------------------------------------------------

\begin{card}
\textbf{\texttt{time.perf\_counter()}}

\begin{lstlisting}[style=pythonstyle]
import time

t0 = time.perf_counter()
result = expensive(n)
print(f"{time.perf_counter() - t0:.3f} s")
\end{lstlisting}

Time a few small $n$, fit the slope of $\log t$ against $\log n$, and
extrapolate before launching a long run.
\end{card}

% ------------------------------------------------------------
\section*{41. Black Boxes: healpy and CAMB}
% ------------------------------------------------------------

\begin{card}
\textbf{healpy} --- Maps on the sphere.

\begin{lstlisting}[style=pythonstyle]
import healpy as hp

nside = 64
print(hp.nside2npix(nside))

alm = hp.map2alm(m, lmax=128)  # map -> a_lm
cl = hp.alm2cl(alm)            # a_lm -> C_l
m2 = hp.alm2map(alm, nside)    # a_lm -> map
cl2 = hp.anafast(m, lmax=128)  # map -> C_l
hp.mollview(m)
\end{lstlisting}

Output (first line): \texttt{49152}, that is $12\,N_{\rm side}^2$ pixels.
\end{card}

\begin{card}
\textbf{CAMB} --- Theory spectra from parameters.

\begin{lstlisting}[style=pythonstyle]
import camb

pars = camb.set_params(H0=67.4, ombh2=0.0224,
                       omch2=0.120, As=2.1e-9,
                       ns=0.965, lmax=2500)
res = camb.get_results(pars)
cl = res.get_cmb_power_spectra(pars,
        CMB_unit="muK", raw_cl=True)["total"][:, 0]
\end{lstlisting}

Six numbers in, $C_\ell^{TT}$ out. In the book this is wrapped in
\texttt{camb\_fiducial.py}.
\end{card}

\begin{card}
\textbf{The book's own helpers} (\texttt{code/common.py})

\begin{itemize}[leftmargin=*,nosep]
  \item \texttt{setup(w, h)}: house plot style and size
  \item \texttt{savefig(fig, "ch06", "name")}: the figure in \texttt{figures/ch06/}
  \item \texttt{save\_numbers("ch06", "name", \{...\})}: numbers as \LaTeX{} macros quoted in the text
  \item \texttt{rng\_for("ch06", "name")}: a seeded generator for that script, the same on every computer
\end{itemize}
\end{card}

\end{multicols}

\end{document}
